% ============================================================================
% Chapter 6: Instanton Derivation
% ============================================================================
% Source: INSTANTON_DERIVATION_CHAIN.md, S5D_TO_SEFF_Q_REDUCTION.md
% Status: FILLED from SoT
% Contamination check: PASSED (no SM terms)
% ============================================================================

\chapter{Instanton Derivation}
\label{ch:instanton}

\source{INSTANTON\_DERIVATION\_CHAIN.md}

\begin{abstract}
How does the metastable junction escape from its metastable configuration? In EDC, the
decay proceeds via \emph{quantum tunneling}---a topological transition through
the energy barrier identified in Chapter~\ref{ch:metastable}. This chapter
develops the instanton formalism: the effective 1D action $S_{\text{eff}}[q]$,
the Euclidean ``bounce'' solution, and the decay rate formula
$\Gamma = A \exp(-\SE/\hbar)$. The key parameters $\kappa$ and $\Lzero/\bdelta$
are deferred to Chapters~\ref{ch:kappa}--\ref{ch:L0delta}.
\end{abstract}

% ----------------------------------------------------------------------------
% CHAPTER SPINE
% ----------------------------------------------------------------------------
\paragraph{Chapter Spine.}
\textbf{Purpose:} Develop the instanton formalism for metastable junction tunneling;
derive the effective 1D action and decay rate formula.
\textbf{Inputs:} (i)~Double-well $V(q)$ from Ch.~\ref{ch:metastable}~[P],
(ii)~barrier height $\VB = 2\Delta m_{np}$ from Ch.~\ref{ch:metastable}~[P],
(iii)~5D action structure~[Def].
\textbf{Outputs:} (i)~Effective action $S_{\text{eff}}[q]$~[Dc],
(ii)~Euclidean action structure $\SE = \kappa \times (L_0/\delta) \times \VB/E_{\text{scale}}$~[Dc],
(iii)~decay rate formula $\Gamma = A \exp(-\SE/\hbar)$~[Der].
\textbf{Dependencies:} Ch.~\ref{ch:metastable} ($V(q)$ structure, $\VB$).
\textbf{Forward links:} $\kappa$ from Ch.~\ref{ch:kappa}; $L_0/\delta$ from Ch.~\ref{ch:L0delta};
assembly in Ch.~\ref{ch:tau_n}.

% ============================================================================
\section{Euclidean Tunneling in EDC}
\label{sec:instanton:setup}

\subsection{Recap: The Tunneling Problem}

From Chapter~\ref{ch:metastable}, the metastable junction occupies a metastable minimum of
the effective potential $V(q)$:
\begin{itemize}
    \item Anchor: $q = 0$ (global minimum, $V = 0$)
    \item Metastable: $q = q_n$ (local minimum, $V = \Delta m_{np}$) \tagP
    \item Barrier: $q = q_B$ (saddle point, $V = 3\Delta m_{np}$) \tagP
\end{itemize}

\textbf{Geometric sector:} The brane-energy contribution
$V_{\mathrm{geom}}(q)$ has been computed along chosen paths
(Appendix~\ref{app:vq_chosen_path}). It is a single-well potential
with minimum at $q = 0$ (Steiner), confirming the geometric restoring
force [Dc]. The full double-well structure---including the secondary
minimum at $q_n$ and the barrier at $q_B$---requires non-geometric terms
that remain [P].

The barrier height from the metastable junction:
\begin{equation}
    \VB = V(q_B) - V(q_n) = 2 \times \Delta m_{np} \approx 2.59 \MeV \tagP
    \label{eq:VB_recall}
\end{equation}

\textbf{The question:} At what rate does the metastable junction tunnel through this barrier?

\subsection{Why Instantons?}

Classical mechanics forbids traversing an energy barrier. Quantum mechanically,
tunneling is exponentially suppressed:
\begin{equation}
    \Gamma \sim \exp\left(-\frac{\text{action}}{\hbar}\right)
\end{equation}

The \emph{instanton} (or ``bounce'') is a classical solution of the Euclidean
equations of motion that mediates the tunneling process. It provides the
action in the exponent.

% ============================================================================
\section{Effective 1D Action}
\label{sec:instanton:action}

\source{S5D\_TO\_SEFF\_Q\_REDUCTION.md}

\subsection{From 5D to 1D: The Reduction Corridor}

The full 5D action of EDC includes bulk, brane, and boundary terms:
\begin{equation}
    S_{\text{total}} = S_{\text{bulk}} + S_{\text{brane}} + S_{\text{GHY}} + S_{\text{junction}}
    \tagDef
\end{equation}

For the tunneling problem, we reduce this to an effective 1D action in the
collective coordinate $q(t)$:

\begin{assumptionbox}[title=Assumptions: 5D$\to$1D Reduction]
    The reduction from 5D dynamics to effective 1D dynamics requires:
    \begin{enumerate}
        \item \textbf{Timescale separation} \tagP: The collective coordinate $q$ is slow
              compared to all other modes (transverse fluctuations, shape modes).
        \item \textbf{Adiabatic relaxation} \tagP: Fast modes relax instantaneously to
              their $q$-dependent equilibria as $q$ evolves.
        \item \textbf{Mode decoupling} \tagP: Cross-couplings between $q$ and transverse
              modes are negligible at leading order.
    \end{enumerate}
    These assumptions are standard for instanton calculations.
    Verification from the full 5D EDC action: \tagOPEN.
\end{assumptionbox}

\begin{derivationbox}[title=5D$\to$1D Reduction (Adiabatic Approximation)]
    Under the above assumptions, the 5D dynamics reduces to effective 1D dynamics.
    The effective action:
    \begin{equation}
        S_{\text{eff}}[q] = \int dt \left[ \frac{1}{2} M(q) \dot{q}^2 - V(q) \right]
        \label{eq:Seff}
    \end{equation}
    where $M(q)$ is the effective mass and $V(q)$ is the effective potential.
    \tagDc
\end{derivationbox}

\subsection{Effective Mass $M(q)$}

The effective mass arises from:
\begin{itemize}
    \item Junction inertia (kinetic energy of moving the junction node)
    \item Brane kinetic energy (deformation of membrane)
    \item Bulk contributions (coupled gravitational modes)
\end{itemize}

\begin{derivationbox}[title=Effective Mass]
    \begin{equation}
        M(q) := \frac{\partial^2 L_{\text{eff}}}{\partial \dot{q}^2}
    \end{equation}
    \textbf{Current status:} In the simplest approximation, we take $M(q) = M$
    (constant). Dimensional analysis within the Nambu--Goto framework gives
    $M \sim \tau R / c^2$ where $\tau$ is the edge tension and $R$ the
    characteristic arm length (Appendix~\ref{app:vq_chosen_path}). The
    numerical coefficient is undetermined and $q$-independence is assumed,
    not verified.
    \tagP
\end{derivationbox}

\begin{edcopenbox}[title=Open: Derive $M(q)$ from 5D]
    The explicit functional form of $M(q)$ is not yet computed from the 5D
    action. The constant-mass approximation is a working ansatz.

    \textbf{Status:} $M(q)$ = [P]. Derivation from $S_{5D}$ = [OPEN].
    \tagOpen
\end{edcopenbox}

\subsection{Effective Potential $V(q)$}

The effective potential is the static energy of the junction configuration:
\begin{equation}
    V(q) := -L_{\text{eff}}(q, \dot{q}=0)
    \tagDc
\end{equation}

The geometric sector $V_{\mathrm{geom}}(q)$ is a single-well potential with
minimum at $q = 0$ (Steiner) [Dc] (Appendix~\ref{app:vq_chosen_path}). The
full $V(q)$ is \emph{postulated} [P] to have double-well structure with minima
at $q = 0$ (anchor) and $q = q_n$ (metastable), and a maximum at $q = q_B$
(barrier). The secondary minimum and barrier require non-geometric terms
($V_{\mathrm{node}}, V_{\mathrm{bulk}}$) that are not yet derived.

% ============================================================================
\section{Euclidean Action and Bounce}
\label{sec:instanton:bounce}

\subsection{Wick Rotation}

For tunneling problems, we continue to imaginary time $\tau = it$:
\begin{equation}
    t \to -i\tau, \quad \dot{q} \to i \frac{dq}{d\tau}
\end{equation}

The Euclidean action is:
\begin{equation}
    \SE[q] = \int d\tau \left[ \frac{1}{2} M \left(\frac{dq}{d\tau}\right)^2 + V(q) \right]
    \label{eq:SE_integral}
    \tagDer
\end{equation}

Note the sign change: $-V(q) \to +V(q)$. In Euclidean signature, the potential
is inverted---barriers become wells and vice versa.

\subsection{The Bounce Solution}

\begin{derivationbox}[title=Bounce Definition]
    The \textbf{bounce} is a classical solution $q_B(\tau)$ of the Euclidean
    equations of motion:
    \begin{equation}
        M \frac{d^2 q}{d\tau^2} = \frac{dV}{dq}
        \label{eq:euclidean_eom}
    \end{equation}
    with boundary conditions:
    \begin{equation}
        q(\tau \to -\infty) = q_n, \quad q(\tau \to +\infty) = q_n
        \label{eq:bounce_bc}
    \end{equation}
    The solution ``bounces'' off the barrier (now a well in inverted potential)
    and returns to the metastable minimum.
    \tagDer
\end{derivationbox}

\textbf{Physical interpretation:} The bounce represents the most probable
tunneling path in the Euclidean path integral. It is the dominant contribution
to the decay amplitude.

\subsection{Euclidean Action of the Bounce}

The action of the bounce solution determines the tunneling rate. For a
general double-well potential, the action is:
\begin{equation}
    \SE = \int_{q_n}^{q_B} dq \sqrt{2 M (V(q) - V(q_n))} \times 2
    \label{eq:SE_wkb}
    \tagDer
\end{equation}

The factor of 2 accounts for the round-trip (out to barrier and back).

\subsection{Parametric Form in EDC}

In EDC, the action takes the form:
\begin{equation}
    \boxed{\frac{\SE}{\hbar} = \kappa \times \frac{\Lzero}{\bdelta}}
    \label{eq:SE_edc}
    \tagP
\end{equation}

where:
\begin{itemize}
    \item $\kappa$ is a topological factor (derived in Chapter~\ref{ch:kappa})
    \item $\Lzero/\bdelta$ is a geometric ratio (analyzed in Chapter~\ref{ch:L0delta})
\end{itemize}

\textbf{Preview:} With $\kappa = 2\pi$ and $\Lzero/\bdelta \approx \pi^2$:
\begin{equation}
    \frac{\SE}{\hbar} = 2\pi \times \pi^2 \approx 2\pi \times 9.87 \approx 62
    \tagP
\end{equation}

% ============================================================================
\section{Decay Rate and Lifetime}
\label{sec:instanton:rate}

\subsection{Tunneling Rate Formula}

The decay rate from the metastable state is:
\begin{equation}
    \Gamma = \Gamma_0 \exp\left(-\frac{\SE}{\hbar}\right)
    \label{eq:Gamma}
    \tagDer
\end{equation}

where $\Gamma_0$ is the prefactor, often called the ``attempt frequency.''

\subsection{Prefactor $A$}

The prefactor arises from fluctuations around the bounce solution. In the
semiclassical approximation:
\begin{equation}
    \Gamma_0 = A \times \frac{\omegaz}{2\pi}
    \label{eq:Gamma0}
    \tagDc
\end{equation}

where:
\begin{itemize}
    \item $\omegaz$ is the oscillation frequency at the metastable minimum
    \item $A$ is a dimensionless factor from the ratio of fluctuation determinants
\end{itemize}

\begin{derivationbox}[title=Prefactor Structure]
    The prefactor $A$ is computed from:
    \begin{equation}
        A = \sqrt{\frac{\SE}{2\pi\hbar}} \times
            \left| \frac{\det'(-\partial_\tau^2 + V''(q_B(\tau)))}
                        {\det(-\partial_\tau^2 + \omegaz^2)} \right|^{-1/2}
    \end{equation}
    where $\det'$ excludes the zero mode associated with time-translation
    invariance of the bounce.
    \tagDer
\end{derivationbox}

\textbf{Order of magnitude:} For smooth potentials, $A = O(1)$.

In practice, $A$ is currently treated as a calibration parameter:
\begin{equation}
    A \approx 0.75 - 0.94 \tagCal
\end{equation}

\begin{edcopenbox}[title=Open: Derive $A$ from fluctuation determinant]
    Computing $A$ requires the explicit bounce solution and potential shape.
    This is deferred to future work.

    \textbf{Status:} $A$ = [Cal]. Fluctuation determinant calculation = [OPEN].
    \tagOpen
\end{edcopenbox}

\subsection{Oscillation Frequency $\omegaz$}

The frequency at the metastable junction minimum:
\begin{equation}
    \omegaz = \sqrt{\frac{V''(q_n)}{M}}
    \tagDc
\end{equation}

\textbf{Dimensional estimate:}
\begin{equation}
    \omegaz \sim \sqrt{\frac{\bsigma}{m_p}} \approx 19 \MeV
    \tagP
\end{equation}

where we use $\bsigma \approx 8.82$ MeV/fm$^2$ and junction mass scale $m_*$.

\subsection{Lifetime Formula}

The metastable junction lifetime is $\taun = 1/\Gamma$:

\begin{resultbox}[title=Metastable Lifetime Formula]
    \begin{equation}
        \boxed{\taun = A \cdot \frac{\hbar}{\omegaz} \cdot \exp\left[\kappa \frac{\Lzero}{\bdelta}\right]}
        \label{eq:tau_n_formula}
    \end{equation}
    \tagDer
\end{resultbox}

This is the central result of the instanton analysis. The exponential
sensitivity to the dimensionless ratio $\Lzero/\bdelta$ explains the large
hierarchy between the junction timescale ($\sim 10^{-24}$ s) and the metastable junction
lifetime ($\sim 10^3$ s).

% ============================================================================
\section{Forward References}
\label{sec:instanton:forward}

The formula \eqref{eq:tau_n_formula} contains parameters that require
further derivation:

\begin{center}
\begin{tabular}{lll}
\toprule
\textbf{Parameter} & \textbf{Status} & \textbf{Derived in} \\
\midrule
$\kappa = 2\pi$ & [P] $\to$ [Dc] & Chapter~\ref{ch:kappa} (homotopy) \\
$\Lzero/\bdelta \approx \pi^2$ & [P] & Chapter~\ref{ch:L0delta} (geometry) \\
$\omegaz \sim 19$ MeV & [P] & — (dimensional estimate) \\
$A \approx 0.9$ & [Cal] & — (from fit) \\
\bottomrule
\end{tabular}
\end{center}

\textbf{Chapter~\ref{ch:kappa}} derives $\kappa = 2\pi$ from the fundamental
group $\pione(\Sone) = \mathbb{Z}$, showing that the topological winding
number is quantized.

\textbf{Chapter~\ref{ch:L0delta}} analyzes the geometric ratio $\Lzero/\bdelta$
and the hypothesis that it equals $\pi^2$.

\textbf{Chapter~\ref{ch:tau_n}} assembles all parameters and computes
$\taun \approx 880\second$.

% ============================================================================
\section{Epistemic Status}
\label{sec:instanton:epistemic}

\begin{center}
\begin{tabular}{llll}
\toprule
\textbf{Object} & \textbf{Meaning} & \textbf{Status} & \textbf{Used in} \\
\midrule
$S_{\text{eff}}[q]$ & 1D effective action & [Dc] & This chapter \\
$M(q)$ & Effective mass & [P] (scaling only) & Eq.~\eqref{eq:Seff} \\
$V_{\mathrm{geom}}(q)$ & Geometric sector & [Dc] (chosen path) & App.~\ref{app:vq_chosen_path} \\
$V(q)$ & Full eff.\ potential & [P] (double-well) & Eq.~\eqref{eq:Seff} \\
$\SE$ & Euclidean action & [Der] (form) & Eq.~\eqref{eq:SE_integral} \\
$\SE/\hbar = \kappa(\Lzero/\bdelta)$ & EDC parametrization & [P] & Eq.~\eqref{eq:SE_edc} \\
$\Gamma = \Gamma_0 e^{-\SE/\hbar}$ & Decay rate & [Der] & Eq.~\eqref{eq:Gamma} \\
$A$ & Prefactor & [Cal] & Eq.~\eqref{eq:Gamma0} \\
$\omegaz$ & Oscillation frequency & [P] & Eq.~\eqref{eq:tau_n_formula} \\
\midrule
$M(q)$ from 5D & Kinetic coefficient & [OPEN] & Future \\
$V(q)$ from 5D & Potential from action & [OPEN] & Future \\
$A$ from det ratio & Fluctuation determinant & [OPEN] & Future \\
\bottomrule
\end{tabular}
\end{center}

% ============================================================================
\section{Falsifiability}
\label{sec:instanton:falsify}

\begin{edcopenbox}[title=Falsifiability Hooks]
    \begin{enumerate}
        \item \textbf{Action form:} If the 5D$\to$1D reduction does not yield
              an action of the form \eqref{eq:Seff}, the instanton framework
              is inapplicable.

        \item \textbf{Timescale separation:} If fast modes do not decouple
              from the collective coordinate, the adiabatic approximation fails.

        \item \textbf{Semiclassical validity:} If $\SE/\hbar \lesssim 1$,
              quantum corrections dominate and the semiclassical formula
              \eqref{eq:Gamma} is invalid. (For metastable junction: $\SE/\hbar \approx 60$,
              safely semiclassical.)

        \item \textbf{Prefactor order:} If $A \ll 1$ or $A \gg 1$ is required
              to match experiment, the model has unnatural fine-tuning.
    \end{enumerate}
\end{edcopenbox}

% ============================================================================
% Observer-side projection
% ============================================================================

\begin{observerbox}
Observer-side projection note: quoted terms are measurement labels for the
3D/4D projection of the 5D objects defined in this chapter; they do not
imply any conventional mechanism.

\begin{itemize}
    \item \MetastableJunction{} $\leftrightarrow$ ``neutron'' (projection label)
    \item \JunctionTransition{} $\leftrightarrow$ measured transition (projection label)
\end{itemize}

What the observer records: the instanton transition rate computed from the
Euclidean action $\SE$ corresponds to the measured lifetime $\taun$ of the
metastable species at the projection level.
\end{observerbox}

% ============================================================================
\section{Summary}
\label{sec:instanton:summary}

\begin{resultbox}
    The metastable junction decay rate is computed via instanton tunneling:
    \begin{enumerate}
        \item The 5D action reduces to effective 1D action $S_{\text{eff}}[q]$
        \item The Euclidean ``bounce'' solution mediates tunneling
        \item The decay rate is $\Gamma = A(\omegaz/2\pi) \exp(-\SE/\hbar)$
        \item The Euclidean action has the form $\SE/\hbar = \kappa(\Lzero/\bdelta)$
    \end{enumerate}

    \textbf{Key result:}
    \begin{equation*}
        \taun = A \cdot \frac{\hbar}{\omegaz} \cdot \exp\left[\kappa \frac{\Lzero}{\bdelta}\right]
    \end{equation*}

    The parameters $\kappa$ and $\Lzero/\bdelta$ are derived in
    Chapters~\ref{ch:kappa}--\ref{ch:L0delta}. Chapter~\ref{ch:tau_n}
    evaluates this formula numerically.
\end{resultbox}

% ----------------------------------------------------------------------------
% BRIDGE TO NEXT CHAPTER
% ----------------------------------------------------------------------------
\paragraph{Bridge to Chapter~\ref{ch:kappa}.}
The instanton formula $\SE/\hbar = \kappa(L_0/\delta)$ contains two
undetermined factors. What fixes $\kappa$? Chapter~\ref{ch:kappa} derives
$\kappa = 2\pi$ from the fundamental group of the circle---the compact
topology of the fifth dimension quantizes the winding number and thereby
fixes the instanton prefactor. This is a topological result, not a fit.
